\subsubsection{Etapa 1.3: servicio operacional heterogéneo}
\label{subsubsec:sp1-heterogeneity}

Dos coaliciones de igual cardinalidad pueden producir valores distintos de servicio $S_k$. E2 conserva la exclusividad, el cierre y las poses estáticas, y asigna a cada pareja robot--carga una contribución calculada con carga útil nominal, batería y distancia. Este índice sirve para comparar disponibilidad operacional. La factibilidad mecánica se evalúa por separado.

\paragraph{Índice operacional de servicio}

El caso didáctico de la Figura~\ref{fig:sp1-compact-e2-scenario} fija $b_i^{\mathrm{res}}=0$, $\delta_{ik}=0$ y $(c_i^{\mathrm{pay}},b_i)\in\{(8,0.5),(8,0.5),(10,0.7),(10,0.5)\}$; por~\eqref{eq:sp2-discounts}--\eqref{eq:sp2-effective-capacity}, $e_{ik}=(4,4,7,5)$. Así, $C_k^A$ suma $8$ y $C_k^B$ suma $12$ frente a $d_k^{\mathrm{srv}}=10$.

% [figure omitido aquí: figura didáctica ya en el cuerpo; ver cuerpo, Sección 6]

Sea $c_i^{\mathrm{pay}}$ la carga útil nominal del robot $i$ en kilogramos, $b_i\in[0, 1]$ su fracción de batería y $b_i^{\mathrm{res}}\in[0, 1)$ la reserva mínima. La distancia inicial robot--carga es $\delta_{ik}$ y $\ell_d>0$ es la escala operacional de distancia. Se definen los factores adimensionales

\begin{equation}
\label{eq:sp2-discounts}
  \psi_i^b=\max\!\left\{0,\frac{b_i-b_i^{\mathrm{res}}}{1-b_i^{\mathrm{res}}}\right\},
  \qquad
  \psi_{ik}^d=\exp\!\left(-\frac{\delta_{ik}}{\ell_d}\right),
  \qquad
  a_{ik}=\psi_i^b\psi_{ik}^d,
\end{equation}

y, con una escala de referencia $c^{\mathrm{ref}}=1\,\mathrm{kg}$, la contribución de servicio

\begin{equation}
\label{eq:sp2-effective-capacity}
  \begin{aligned}
    e_{ik}&=\frac{c_i^{\mathrm{pay}}}{c^{\mathrm{ref}}}a_{ik},
    & S_k(x)&=\sum_{i=1}^{N}e_{ik}x_{ik},\\
    S_k(x)&\geq d_k^{\mathrm{srv}}
    &&\Longleftrightarrow \text{ umbral operacional cubierto}.
  \end{aligned}
\end{equation}

El factor $a_{ik}$ representa disponibilidad y $e_{ik}$ expresa servicio normalizado, no masa. La elección $c^{\mathrm{ref}}=1\,\mathrm{kg}$ fija la escala, mientras $\psi_{ik}^d$ penaliza el desplazamiento inicial sin modificar la carga útil nominal. Los experimentos suponen compatibilidad unitaria y omiten restricciones independientes de soporte y balance energético.

\paragraph{Problemas formales de referencia}

Cobertura parcial y completitud son estimandos distintos. Dos MILP globales usan $x_{ik}\in\{0,1\}$ y servicio cubierto $s_k\in[0,d_k^{\mathrm{srv}}]$, con exclusividad por robot.

La referencia de cobertura normalizada es

\begin{equation}
\label{eq:sp2-capacity-oracle}
\begin{aligned}
  \max_{x,s}\quad
  & \sum_{k=1}^{K}\frac{s_k}{d_k^{\mathrm{srv}}}
    -10^{-5}\sum_{i,k}\delta_{ik}x_{ik} \\
  \text{sujeto a}\quad
  & \sum_k x_{ik}\leq1, && \forall i,\\
  & 0\leq s_k\leq d_k^{\mathrm{srv}},\quad s_k\leq\sum_i e_{ik}x_{ik}, && \forall k,\\
  & x_{ik}\in\{0, 1\}. &&
\end{aligned}
\end{equation}

En~\eqref{eq:sp2-capacity-oracle}, el término de $10^{-5}$ puntos/m resuelve empates después de maximizar la suma de razones cubiertas. La diferencia respecto a este objetivo mide cobertura normalizada y no constituye una cota para completitud u otras métricas.

La referencia de puntuación añade $z_k\in\{0, 1\}$, activo solo si se cumplen demanda y cardinalidad $n_k$, y usa los pesos del experimento:

\begin{equation}
\label{eq:sp2-score-oracle}
\begin{aligned}
  \max_{x,z,s}\quad
  & 5\sum_k V_kz_k
    +100\sum_k\frac{s_k}{d_k^{\mathrm{srv}}}
    -5\times10^{-4}\sum_{i,k}\delta_{ik}x_{ik}
    -10^{-5}\sum_{i,k}E_{ik}^{\mathrm{trav}}x_{ik} \\
  \text{sujeto a}\quad
  & \sum_kx_{ik}\leq1, && \forall i,\\
  & \sum_i e_{ik}x_{ik}\geq d_k^{\mathrm{srv}}z_k, && \forall k,\\
  & \sum_i x_{ik}\geq n_kz_k, && \forall k,\\
  & 0\leq s_k\leq d_k^{\mathrm{srv}},\quad s_k\leq\sum_i e_{ik}x_{ik}, && \forall k,\\
  & x_{ik},z_k\in\{0, 1\}. &&
\end{aligned}
\end{equation}

En~\eqref{eq:sp2-score-oracle}, $V_k$ es adimensional y $E_{ik}^{\mathrm{trav}}$ se mide en Wh, de modo que los coeficientes tienen unidades inversas. Los dos MILP reciben la matriz completa $[e_{ik}]$, las demandas y los costes; se consideran exactos únicamente en las instancias cuyo optimizador certifica brecha nula.

\paragraph{Métodos de comparación}

Los dos MILP fijan referencias globales para cobertura y completitud. Frente a ellos se evalúan asignadores con información agregada, reglas locales o vecinales y modelos de imitación. El Húngaro expandido y CBBA-capacidad modifican los problemas estudiados en sus fuentes, por lo que sus garantías originales no se transfieren a esta comparación \parencite{kuhn1955Hungarian,choiBrunetHow2009CBBA}. En la Tabla~\ref{tab:sp1-compact-e2-methods}, $S$ denota puestos, $d$ rasgos y $P$ pesos.

% [table omitido aquí: tabla de métodos ya en el cuerpo; ver cuerpo, Sección 6]

\paragraph{Juego de capacidad y alineación marginal}

\begin{contribucion}{puntuación de servicio y potencial marginal de E2}
\estadoaporte{demostrado para preferencias continuas y parámetros congelados.}
Durante una decisión se fija la matriz $E=[e_{ik}]$ para separar el efecto de la heterogeneidad del cambio de estado. La variable $\rho_{ik}\in[0,1]$ expresa la preferencia continua del robot $i$ por la carga $k$ y satisface $\sum_k\rho_{ik}\leq1$. Con la presión saturada $\sigma(r)=[1-r]_+$ aplicada a $r_k=S_k(\rho)/d_k^{\mathrm{srv}}$ y un coste normalizado $g_{ik}$, las puntuaciones son

\begin{equation}
\label{eq:sp2-payoffs}
  f_{ik}^{\mathrm{plain}}
  =V_k\left[1-\frac{S_k(\rho)}{d_k^{\mathrm{srv}}}\right]_+-g_{ik},
  \qquad
  f_{ik}^{\mathrm{marg}}
  =\frac{e_{ik}}{d_k^{\mathrm{srv}}}V_k\left[1-\frac{S_k(\rho)}{d_k^{\mathrm{srv}}}\right]_+-g_{ik}.
\end{equation}

El factor $e_{ik}/d_k^{\mathrm{srv}}$ distingue la fracción de demanda que aporta cada robot. Sin él, $f_{ik}^{\mathrm{plain}}$ asigna la misma presión a contribuciones diferentes.

\begin{proposicion}[Alineación marginal del servicio operacional]
\label{prop:sp2-marginal-alignment-detail}
En una región donde $S_k(\rho)<d_k^{\mathrm{srv}}$, con $E$, $d_k^{\mathrm{srv}}$, $V_k$ y $g_{ik}$ fijos, el campo marginal de~\eqref{eq:sp2-payoffs} es el gradiente del potencial
\begin{equation}
\label{eq:sp2-potential}
  \Phi(\rho)=
  \sum_k V_k\int_0^{S_k(\rho)/d_k^{\mathrm{srv}}}[1-s]_+\,\mathrm ds
  -\sum_{i,k}g_{ik}\rho_{ik}.
\end{equation}
El campo plano no es, en general, integrable como potencial exacto cuando existen $i,j,k$ con $e_{ik}\neq e_{jk}$. La identidad se establece para preferencias continuas, parámetros fijos y una región no saturada.
\end{proposicion}
\end{contribucion}

Por la regla de la cadena, $\partial\Phi/\partial\rho_{ik}=f_{ik}^{\mathrm{marg}}$; al retirar el factor $e_{ik}/d_k^{\mathrm{srv}}$, las derivadas cruzadas dejan de coincidir cuando las contribuciones son heterogéneas. La frontera saturada se trata en el Anexo~\ref{app:sp2-marginal-proof}.

El contrato distribuido requiere que cada robot conozca $e_{ik}$ y $(d_k^{\mathrm{srv}},V_k)$, además de estimar $S_k$ con mensajes vecinales. En la campaña, las reglas recibieron el déficit global salvo la variante primal--dual de radio limitado; las variantes Smith se implementaron como ordenaciones secuenciales y no como integraciones de la EDO.

\paragraph{Regulación del déficit de capacidad}

La salida regulada en E2 es el déficit relativo de servicio:
\begin{equation}
 z_{2,k}(\rho)=
 \left[1-\frac{S_k(\rho)}{d_k^{\mathrm{srv}}}\right]_+,
 \qquad z_{2,k}^{\mathrm{ref}}=0
 \label{eq:sp2-strategic-regulation}
\end{equation}
para cada carga de~\eqref{eq:sp2-strategic-regulation}. Si no pueden cubrirse todas, el optimizador selecciona cuáles regulan $z_{2,k}=0$ y las demás se liberan.

La realimentación mínima del robot $i$ está formada por su contribución $e_{ik}$, su coste $g_{ik}$, los parámetros públicos $(d_k^{\mathrm{srv}},V_k)$ y una estimación $\widehat S_{ik}$ del servicio agregado. Con esos datos calcula
\begin{equation}
 \widehat f_{ik}^{\mathrm{marg}}
 =\frac{e_{ik}}{d_k^{\mathrm{srv}}}V_k
  \left[1-\frac{\widehat S_{ik}}{d_k^{\mathrm{srv}}}\right]_+-g_{ik},
 \qquad
 \rho_i^{m+1}=P_{\Delta_i}\!\left(
 \rho_i^m+\Delta t\,F_i(\widehat f_i^m)\right),
 \label{eq:sp2-feedback-candidate}
\end{equation}
En la realimentación de~\eqref{eq:sp2-feedback-candidate}, $F_i$ puede ser Smith, BNN, replicator o gradiente proyectado. La Proposición~\ref{prop:sp2-marginal-alignment-detail} establece la identidad $\partial\Phi/\partial\rho_{ik}=f_{ik}^{\mathrm{marg}}$ para agregados exactos y $E$ fija. El cierre entero y la estabilidad muestreada requieren resultados adicionales.

La campaña ejecutó una ordenación secuencial con $\widehat S_{ik}=S_k$ en casi todas las reglas y mantuvo fijas las posiciones, la batería y $E$. Este contraste aísla la puntuación; la ley continua y el observador vecinal no forman parte del estimando.
\paragraph{Diseño experimental}

Las semillas $3100$--$3139$ producen $\SPTwoWorlds$ mundos pareados por método y $\SPTwoComparisonRows$ observaciones método--mundo. El endpoint principal es la completitud; la Tabla~\ref{tab:sp2-design} resume los escenarios, las métricas auxiliares y la inferencia. La ablación plano--marginal reutiliza los mundos para cuatro reglas y las dos referencias, con $\SPTwoAblationRows$ comprobaciones válidas.

\begin{table}[!htbp]
  \centering
  \caption{Factores y estimandos de la campaña E2.}
  \label{tab:sp2-design}
  \viutablefont
  \begin{tabularx}{0.94\textwidth}{>{\raggedright\arraybackslash}p{2.5cm} Y}
    \toprule
    Elemento & Especificación \\
    \midrule
    Escenarios & Ligero mixto, capacidad balanceada, capacidad pesada, batería limitada y Monte Carlo. \\
    Métricas & Completitud y cobertura; brechas frente a los MILP; servicio incompleto, alineación, distancia, energía, mensajes y CPU. \\
    Inferencia & Friedman/Kendall y Wilcoxon por mundo, medias pareadas con \mbox{bootstrap} y corrección de Holm; contrastes preespecificados para puntuación marginal y primal--dual frente a voraz. \\
    \bottomrule
  \end{tabularx}
  \viusource{elaboración propia a partir del protocolo experimental de E2}
\end{table}

Los datos publicados permiten recalcular el postproceso y reconstruir cada tabla y figura. La ausencia del generador de instancias impide repetir la creación de mundos y limita la reproducibilidad de la simulación.

\paragraph{Resultados}

Frente a la referencia de cobertura ($\SPTwoCoverageCoverage$ de cobertura media, $\SPTwoCoverageSuccess$ de completitud), el oráculo de puntuación invierte el reparto: cobertura $\SPTwoScoreCoverage$, completitud $\SPTwoScoreSuccess$. La diferencia cuantifica el intercambio entre repartir servicio parcial y cerrar cargas en los dos objetivos MILP (Tabla~\ref{tab:sp2-main-results}).

\begin{table}[!htbp]
  \centering
  \caption[E2 en 1560 mundos pareados. Diferencia de cobertura frente al oráculo de cobertura.]{E2 en $\SPTwoWorlds$ mundos pareados. ``Dif. cobertura'' se define en~\eqref{eq:sp2-capacity-oracle} y no constituye una brecha universal.}
  \label{tab:sp2-main-results}
  \resizebox{\textwidth}{!}{\begin{tabular}{lrrrrr}
\toprule
Método & Cobertura & Completitud & Brecha & Dif. cobertura & CPU [ms] \\
\midrule
Oráculo de puntuación & 0.514 & 0.487 & 0.000 & 0.029 & 40.84 \\
Referencia de cobertura & 0.529 & 0.227 & 0.049 & 0.000 & 11.31 \\
Puntuación neuronal & 0.507 & 0.263 & 0.109 & 0.050 & 1.01 \\
Imitación lineal & 0.514 & 0.218 & 0.118 & 0.038 & 0.98 \\
Voraz de servicio & 0.510 & 0.219 & 0.276 & 0.044 & 0.30 \\
Puntuación inspirada en replicator & 0.525 & 0.139 & 0.297 & 0.026 & 0.36 \\
Puntuación inspirada en Smith & 0.521 & 0.103 & 0.343 & 0.033 & 0.35 \\
\bottomrule
\end{tabular}
}
  \viusource{elaboración propia a partir de los datos procesados de E2}
\end{table}

Entre los métodos aproximados, la red neuronal obtuvo la menor brecha de puntuación, $\SPTwoNeuralGap$, con completitud $\SPTwoNeuralSuccess$ y $\SPTwoNeuralRuntime\,\mathrm{ms}$ de CPU media. Smith fue más barata, $\SPTwoSmithRuntime\,\mathrm{ms}$, pero alcanzó una brecha de $\SPTwoSmithGap$ y completitud $\SPTwoSmithSuccess$. La regla interpretable reduce el coste computacional, aunque no encabeza la calidad de asignación.

La ablación de la Tabla~\ref{tab:sp2-ablation} evalúa el factor marginal dentro del asignador secuencial. Frente a la puntuación plana, la corrección elevó la completitud de $\SPTwoPlainSuccess$ a $\SPTwoMarginalSuccess$ y redujo la brecha de $\SPTwoPlainGap$ a $\SPTwoMarginalGap$; el efecto pareado sobre la brecha fue $\SPTwoGapEffect$ (IC del 95\,\% $[\SPTwoGapCILow,\SPTwoGapCIHigh]$, $p_{\mathrm{Holm}}=\SPTwoGapPHolm$). El cambio en completitud fue $\SPTwoSuccessEffect$ (IC $[\SPTwoSuccessCILow,\SPTwoSuccessCIHigh]$) y el del servicio retenido en cargas incompletas, $\SPTwoIncompleteEffect$. La campaña no integró la EDO Smith.

\begin{table}[!htbp]
  \centering
  \caption{Ablación pareada de las puntuaciones plana y marginal en E2.}
  \label{tab:sp2-ablation}
  \resizebox{0.82\textwidth}{!}{\begin{tabular}{lrrrr}
\toprule
Variante & Completitud & Brecha & Serv. incompleto & Alineación \\
\midrule
Puntuación plana & 0.051 & 0.372 & 0.946 & 0.049 \\
Puntuación marginal & 0.135 & 0.157 & 0.848 & 0.137 \\
\bottomrule
\end{tabular}
}
  \viusource{elaboración propia a partir de los datos procesados de E2}
\end{table}

Para la variante primal--dual ejecutada, la diferencia media de completitud frente al voraz fue $\SPTwoPDEffect$ (IC del 95\,\% $[\SPTwoPDCILow,\SPTwoPDCIHigh]$, $p_{\mathrm{Holm}}=\SPTwoPDPHolm$). Sus selecciones dispersaron capacidad entre cargas sin reflejar la prioridad todo-o-nada. Esta observación describe la parametrización ensayada; la corrección marginal, en cambio, sí mejoró la regla secuencial.

Dos amenazas dominan la validez externa. El índice mezcla disponibilidad y distancia sin un modelo mecánico, de modo que una asignación operacionalmente cubierta puede ser físicamente inviable. El uso generalizado del déficit global impide atribuir los resultados a una implementación vecinal. Los rangos reportados corresponden solo a los mundos evaluados.

\paragraph{Síntesis}

La heterogeneidad rompe la equivalencia entre cardinalidad y cobertura, y el objetivo agregado puede repartir capacidad sin completar cargas. Con $E$ fija, el factor $e_{ik}/d_k^{\mathrm{srv}}$ alinea el campo con el gradiente de~\eqref{eq:sp2-potential}; en la ablación también concentró mejor el servicio. Entre los métodos aproximados, la red neuronal obtuvo la menor brecha y la variante primal--dual quedó por debajo del voraz en completitud.

El umbral $S_k\geq d_k^{\mathrm{srv}}$ puede cumplirse aunque los robots carezcan de puntos de contacto o de brazos de palanca capaces de equilibrar el torque. Ese contraejemplo motiva el certificado geométrico de \emph{wrench} de E3.
